% GENERATED FILE. Edit canonical lessons, then run npm run generate:books.
% canonical-source-hash: fnv1a64:f2ea3656a1840c75
% canonical-lessons: PA-C280-bhan, PA-C280-bhanauna, PA-C280-bachna, PA-C280-puggna, PA-C280-baki, PA-R280-review-274-280, PA-R280-review-final-280

\chapter{Small Change}
\label{ch:pa-a2-small-change}
\hlchaptermodality{280}

\begin{chapteropening}
\textbf{By the end of this chapter:} \emph{I can ask for change and say what money is left.}

Complete the last lesson of chapter 280: I can ask for change and say what money is left.
\end{chapteropening}

\section[\texorpdfstring{\pa{ਭਾਨ} (bhān)}{ਭਾਨ (bhān)}]{\pa{ਭਾਨ} (bhān) — small change; coins}

\label{lesson:PA-C280-bhan}



\begin{quote}
Before the new one: say the Punjabi for to exchange, then the Punjabi for a complaint.
\end{quote}

\subsection*{You'll want to know: \pa{ਭਾਨ}}
\textbf{\pa{ਭਾਨ}} — \emph{bhān} — \textquotedblleft{}small change; coins\textquotedblright{}.

\begin{quote}
\pa{ਕੀ} \pa{ਤੁਹਾਡੇ} \pa{ਕੋਲ} \pa{ਭਾਨ} \pa{ਹੈ}?
\end{quote}

\begin{quote}
\emph{kī tuhāḍe kol bhān hai?}
\end{quote}

\begin{quote}
Do you have small change?
\end{quote}

A feminine noun.

\begin{cousinweb}
\textbf{\pa{ਭਾਨ}} (\emph{bhān}) is inherited from Sanskrit \emph{bhājana}, \textquotedblleft{}dividing\textquotedblright{}; CDIAL lists Punjabi \emph{bhāṇ}, \textquotedblleft{}small change for a rupee\textquotedblright{}, under it (CDIAL 9437).
\end{cousinweb}

\begin{grammarlens}[title={What you've built}]
One more word for this chapter.
\end{grammarlens}

\subsection*{Your turn}
\begin{itemize}
\raggedright
  \item \emph{Say it:} \emph{bhān}
  \item \emph{Say it:} \emph{bhān}, once more
  \item \emph{Say it:} \emph{kī tuhāḍe kol bhān hai?} — Do you have small change?
  \item \emph{Recall:} say the Punjabi for to exchange, then the Punjabi for a complaint, then say \emph{bhān} again
\end{itemize}

\begin{tcolorbox}[breakable,colback=teal!4,colframe=teal!35!black,title={Before you move on}]
What does \pa{ਭਾਨ} mean? (\textquotedblleft{}Small change; coins\textquotedblright{}.) Say it once more.
\end{tcolorbox}
\section[\texorpdfstring{\pa{ਭਨਾਉਣਾ} (bhanāuṇā)}{ਭਨਾਉਣਾ (bhanāuṇā)}]{\pa{ਭਨਾਉਣਾ} (bhanāuṇā) — to break a note into change}

\label{lesson:PA-C280-bhanauna}



\begin{quote}
Before the new one: say the Punjabi for a complaint, then the Punjabi for small change.
\end{quote}

\subsection*{You'll want to know: \pa{ਭਨਾਉਣਾ}}
\textbf{\pa{ਭਨਾਉਣਾ}} — \emph{bhanāuṇā} — \textquotedblleft{}to break a note into change\textquotedblright{}.

\begin{quote}
\pa{ਮੈਂ} \pa{ਬੈਂਕ} \pa{ਵਿੱਚ} \pa{ਸੌ} \pa{ਦਾ} \pa{ਨੋਟ} \pa{ਭਨਾਉਂਦਾ} \pa{ਹਾਂ।}
\end{quote}

\begin{quote}
\emph{maiṁ baiṅk vich sau dā noṭ bhanāundā hāṁ.}
\end{quote}

\begin{quote}
I break a hundred note into change at the bank.
\end{quote}

Past (I …): \pa{ਮੈਂ} \pa{ਨੋਟ} \pa{ਭਨਾਇਆ}, \emph{maiṁ noṭ bhanāiā}.

\begin{cousinweb}
\textbf{\pa{ਭਨਾਉਣਾ}} (\emph{bhanāuṇā}) is inherited from Sanskrit \emph{bhajyate}, \textquotedblleft{}is broken\textquotedblright{}; CDIAL lists Punjabi \emph{bhanāuṇā}, \textquotedblleft{}to change (money)\textquotedblright{}, under it (CDIAL 9361).
\end{cousinweb}

\begin{grammarlens}[title={What you've built}]
One more word for this chapter.
\end{grammarlens}

\subsection*{Your turn}
\begin{itemize}
\raggedright
  \item \emph{Say it:} \emph{bhanāuṇā}
  \item \emph{Say it:} \emph{bhanāuṇā}, once more
  \item \emph{Say it:} \emph{maiṁ baiṅk vich sau dā noṭ bhanāundā hāṁ.} — I break a hundred note into change at the bank.
  \item \emph{Recall:} say the Punjabi for a complaint, then the Punjabi for small change, then say \emph{bhanāuṇā} again
  \item \emph{Write it:} copy \textbf{\pa{ਭਨਾਉਣਾ}} once, checking the whole word
\end{itemize}

\begin{tcolorbox}[breakable,colback=teal!4,colframe=teal!35!black,title={Before you move on}]
What does \pa{ਭਨਾਉਣਾ} mean? (\textquotedblleft{}To break a note into change\textquotedblright{}.) Say it once more.
\end{tcolorbox}
\section[\texorpdfstring{\pa{ਬਚਣਾ} (bachṇā)}{ਬਚਣਾ (bachṇā)}]{\pa{ਬਚਣਾ} (bachṇā) — to be left over; to remain}

\label{lesson:PA-C280-bachna}



\begin{quote}
Before the new one: say the Punjabi for small change, then the Punjabi for to break a note into change.
\end{quote}

\subsection*{You'll want to know: \pa{ਬਚਣਾ}}
\textbf{\pa{ਬਚਣਾ}} — \emph{bachṇā} — \textquotedblleft{}to be left over; to remain\textquotedblright{}.

\begin{quote}
\pa{ਸੌਦੇ} \pa{ਤੋਂ} \pa{ਬਾਅਦ} \pa{ਮੇਰੇ} \pa{ਕੋਲ} \pa{ਕੁਝ} \pa{ਪੈਸੇ} \pa{ਬਚਦੇ} \pa{ਹਨ।}
\end{quote}

\begin{quote}
\emph{saude ton bāad mere kol kujh paise bachde han.}
\end{quote}

\begin{quote}
After the shopping I have some money left.
\end{quote}

Past (it …): \pa{ਉਹ} \pa{ਰੋਟੀ} \pa{ਬਚੀ}, \emph{oh roṭī bachī}.

\begin{cousinweb}
\textbf{\pa{ਬਚਣਾ}} (\emph{bachṇā}) goes back to Sanskrit \emph{vañcati}, \textquotedblleft{}moves; escapes from\textquotedblright{}; CDIAL lists Punjabi \emph{bacṇā}, \textquotedblleft{}to escape\textquotedblright{}, under it, as a word taken from Hindi \emph{bacnā} (CDIAL 11208).
\end{cousinweb}

\begin{grammarlens}[title={What you've built}]
One more word for this chapter.
\end{grammarlens}

\subsection*{Your turn}
\begin{itemize}
\raggedright
  \item \emph{Say it:} \emph{bachṇā}
  \item \emph{Say it:} \emph{bachṇā}, once more
  \item \emph{Say it:} \emph{saude ton bāad mere kol kujh paise bachde han.} — After the shopping I have some money left.
  \item \emph{Recall:} say the Punjabi for small change, then the Punjabi for to break a note into change, then say \emph{bachṇā} again
  \item \emph{Write it:} copy \textbf{\pa{ਬਚਣਾ}} once, checking the whole word
\end{itemize}

\begin{tcolorbox}[breakable,colback=teal!4,colframe=teal!35!black,title={Before you move on}]
What does \pa{ਬਚਣਾ} mean? (\textquotedblleft{}To be left over; to remain\textquotedblright{}.) Say it once more.
\end{tcolorbox}
\section[\texorpdfstring{\pa{ਪੁੱਗਣਾ} (puggṇā)}{ਪੁੱਗਣਾ (puggṇā)}]{\pa{ਪੁੱਗਣਾ} (puggṇā) — to be affordable; to be worth it}

\label{lesson:PA-C280-puggna}



\begin{quote}
Before the new one: say the Punjabi for to break a note into change, then the Punjabi for to be left over.
\end{quote}

\subsection*{You'll want to know: \pa{ਪੁੱਗਣਾ}}
\textbf{\pa{ਪੁੱਗਣਾ}} — \emph{puggṇā} — \textquotedblleft{}to be affordable; to be worth it\textquotedblright{}.

\begin{quote}
\pa{ਮਹਿੰਗਾ} \pa{ਕਿਰਾਇਆ} \pa{ਮੈਨੂੰ} \pa{ਨਹੀਂ} \pa{ਪੁੱਗਦਾ} \pa{ਹੈ।}
\end{quote}

\begin{quote}
\emph{mahiṅgā kirāiā mainū̃ nahīṁ puggdā hai.}
\end{quote}

\begin{quote}
I cannot afford an expensive fare.
\end{quote}

Past (it …): \pa{ਉਹ} \pa{ਸੌਦਾ} \pa{ਮੈਨੂੰ} \pa{ਪੁੱਗਿਆ}, \emph{oh saudā mainū̃ puggiā}.

\begin{cousinweb}
\textbf{\pa{ਪੁੱਗਣਾ}} (\emph{puggṇā}) is inherited from Sanskrit \emph{pūryate}, \textquotedblleft{}is filled\textquotedblright{}; CDIAL lists Punjabi \emph{puggṇā}, \textquotedblleft{}to be completed, reach\textquotedblright{}, under it (CDIAL 8342). The sense \textquotedblleft{}to be affordable\textquotedblright{} is not given there.
\end{cousinweb}

\begin{grammarlens}[title={What you've built}]
One more word for this chapter.
\end{grammarlens}

\subsection*{Your turn}
\begin{itemize}
\raggedright
  \item \emph{Say it:} \emph{puggṇā}
  \item \emph{Say it:} \emph{puggṇā}, once more
  \item \emph{Say it:} \emph{mahiṅgā kirāiā mainū̃ nahīṁ puggdā hai.} — I cannot afford an expensive fare.
  \item \emph{Recall:} say the Punjabi for to break a note into change, then the Punjabi for to be left over, then say \emph{puggṇā} again
  \item \emph{Write it:} copy \textbf{\pa{ਪੁੱਗਣਾ}} once, checking the whole word
\end{itemize}

\begin{tcolorbox}[breakable,colback=teal!4,colframe=teal!35!black,title={Before you move on}]
What does \pa{ਪੁੱਗਣਾ} mean? (\textquotedblleft{}To be affordable; to be worth it\textquotedblright{}.) Say it once more.
\end{tcolorbox}
\section[\texorpdfstring{\pa{ਬਾਕੀ} (bākī)}{ਬਾਕੀ (bākī)}]{\pa{ਬਾਕੀ} (bākī) — the rest; remaining}

\label{lesson:PA-C280-baki}



\begin{quote}
Before the new one: say the Punjabi for to be left over, then the Punjabi for to be affordable.
\end{quote}

\subsection*{You'll want to know: \pa{ਬਾਕੀ}}
\textbf{\pa{ਬਾਕੀ}} — \emph{bākī} — \textquotedblleft{}the rest; remaining\textquotedblright{}.

\begin{quote}
\pa{ਬਾਕੀ} \pa{ਪੈਸੇ} \pa{ਮੇਰੀ} \pa{ਜੇਬ} \pa{ਵਿੱਚ} \pa{ਹਨ।}
\end{quote}

\begin{quote}
\emph{bākī paise merī jeb vich han.}
\end{quote}

\begin{quote}
The rest of the money is in my pocket.
\end{quote}

\begin{cousinweb}
\textbf{\pa{ਬਾਕੀ}} (\emph{bākī}) is a loan, usually traced through Persian to Arabic \emph{bāqī}, \textquotedblleft{}remaining\textquotedblright{}.
\end{cousinweb}

\begin{grammarlens}[title={What you've built}]
One more word for this chapter.
\end{grammarlens}

\subsection*{Your turn}
\begin{itemize}
\raggedright
  \item \emph{Say it:} \emph{bākī}
  \item \emph{Say it:} \emph{bākī}, once more
  \item \emph{Say it:} \emph{bākī paise merī jeb vich han.} — The rest of the money is in my pocket.
  \item \emph{Recall:} say the Punjabi for to be left over, then the Punjabi for to be affordable, then say \emph{bākī} again
  \item \emph{Write it:} copy \textbf{\pa{ਬਾਕੀ}} once, checking the whole word
\end{itemize}

\begin{tcolorbox}[breakable,colback=teal!4,colframe=teal!35!black,title={Before you move on}]
What does \pa{ਬਾਕੀ} mean? (\textquotedblleft{}The rest; remaining\textquotedblright{}.) Say it once more.
\end{tcolorbox}
\section[(dialogue)]{Looking back: from Yards and Metres to Small Change}

\label{lesson:PA-R280-review-274-280}



\begin{quote}
Say the Punjabi for to measure and a yard.
\end{quote}

\subsection*{Your turn}
Say these aloud:

\begin{itemize}
\raggedright
  \item to measure, a yard, a metre, a size, a pair
  \item loose, tight, to suit, a dozen, a bag
  \item a price, a rate, to lower, a purchase, free of charge
  \item a customer, to choose, to test, a thing, goods
  \item a reckoning, to load, to be sold, a packet, to add up
  \item faulty, fake, genuine, to exchange, a complaint
  \item small change, to break a note into change, to be left over, to be affordable, the rest
\end{itemize}

\begin{tcolorbox}[breakable,colback=teal!4,colframe=teal!35!black,title={Before you move on}]
To measure? (\textbf{\pa{ਮਿਣਨਾ}}, \emph{miṇnā}.) A yard? (\textbf{\pa{ਗਜ਼}}, \emph{gaz}.) A metre? (\textbf{\pa{ਮੀਟਰ}}, \emph{mīṭar}.) A size? (\textbf{\pa{ਨਾਪ}}, \emph{nāp}.) A pair? (\textbf{\pa{ਜੋੜਾ}}, \emph{joṛā}.) Loose? (\textbf{\pa{ਢਿੱਲਾ}}, \emph{ḍhillā}.) Tight? (\textbf{\pa{ਤੰਗ}}, \emph{taṅg}.) To suit? (\textbf{\pa{ਫੱਬਣਾ}}, \emph{phabbṇā}.) A dozen? (\textbf{\pa{ਦਰਜਨ}}, \emph{darjan}.) A bag? (\textbf{\pa{ਥੈਲਾ}}, \emph{thailā}.) A price? (\textbf{\pa{ਮੁੱਲ}}, \emph{mull}.) A rate? (\textbf{\pa{ਭਾਅ}}, \emph{bhāa}.) To lower? (\textbf{\pa{ਘਟਾਉਣਾ}}, \emph{ghaṭāuṇā}.) A purchase? (\textbf{\pa{ਸੌਦਾ}}, \emph{saudā}.) Free of charge? (\textbf{\pa{ਮੁਫ਼ਤ}}, \emph{mufat}.) A customer? (\textbf{\pa{ਗਾਹਕ}}, \emph{gāhak}.) To choose? (\textbf{\pa{ਚੁਣਨਾ}}, \emph{chuṇnā}.) To test? (\textbf{\pa{ਪਰਖਣਾ}}, \emph{parakhṇā}.) A thing? (\textbf{\pa{ਚੀਜ਼}}, \emph{chīz}.) Goods? (\textbf{\pa{ਸਾਮਾਨ}}, \emph{sāmān}.) A reckoning? (\textbf{\pa{ਹਿਸਾਬ}}, \emph{hisāb}.) To load? (\textbf{\pa{ਲੱਦਣਾ}}, \emph{laddṇā}.) To be sold? (\textbf{\pa{ਵਿਕਣਾ}}, \emph{vikṇā}.) A packet? (\textbf{\pa{ਪੈਕਟ}}, \emph{paikaṭ}.) To add up? (\textbf{\pa{ਜੋੜਨਾ}}, \emph{joṛnā}.) Faulty? (\textbf{\pa{ਖ਼ਰਾਬ}}, \emph{kharāb}.) Fake? (\textbf{\pa{ਨਕਲੀ}}, \emph{naklī}.) Genuine? (\textbf{\pa{ਅਸਲੀ}}, \emph{aslī}.) To exchange? (\textbf{\pa{ਵਟਾਉਣਾ}}, \emph{vaṭāuṇā}.) A complaint? (\textbf{\pa{ਸ਼ਿਕਾਇਤ}}, \emph{shikāit}.) Small change? (\textbf{\pa{ਭਾਨ}}, \emph{bhān}.) To break a note into change? (\textbf{\pa{ਭਨਾਉਣਾ}}, \emph{bhanāuṇā}.) To be left over? (\textbf{\pa{ਬਚਣਾ}}, \emph{bachṇā}.) To be affordable? (\textbf{\pa{ਪੁੱਗਣਾ}}, \emph{puggṇā}.) The rest? (\textbf{\pa{ਬਾਕੀ}}, \emph{bākī}.)
\end{tcolorbox}
\section[(dialogue)]{Small Change — once more}

\label{lesson:PA-R280-review-final-280}



\begin{quote}
Say the Punjabi for small change and to break a note into change.
\end{quote}

\subsection*{Your turn}
\begin{itemize}
\raggedright
  \item \emph{Say it:} small change, to break a note into change, to be left over, to be affordable, the rest
\end{itemize}

\begin{tcolorbox}[breakable,colback=teal!4,colframe=teal!35!black,title={Before you move on}]
Small change? (\textbf{\pa{ਭਾਨ}}, \emph{bhān}.) To break a note into change? (\textbf{\pa{ਭਨਾਉਣਾ}}, \emph{bhanāuṇā}.) To be left over? (\textbf{\pa{ਬਚਣਾ}}, \emph{bachṇā}.) To be affordable? (\textbf{\pa{ਪੁੱਗਣਾ}}, \emph{puggṇā}.) The rest? (\textbf{\pa{ਬਾਕੀ}}, \emph{bākī}.)
\end{tcolorbox}
